\section{Details on Model Implementation}\label{app:model-impl}

\subsection{State space in VFI Toolkit terms}
In VFI Toolkit notation the problem has

\begin{itemize}
\item \textbf{Permanent types} $N_i = 4$, \verb|Names_i = {NCL,NCH,CL,CH}|: non-college / college $\times$ low / high productivity (LW2015 Fig.~1). The first two PTypes are assigned to the blue-collar occupational pension scheme, the second two to white-collar.
\item \textbf{Decision variable} \verb|n_d = 5|: an action code dispatching joint labor / claim decisions; see Section~\ref{sec:histidx} for the enumeration.
\item \textbf{Standard endogenous states} \verb|n_a = [n_k, n_hist] = [300, 1237]|: assets $k\in[0,50\text{BA}]$ on 300 grid points, plus a single history-index $\text{hist\_idx}$ that encodes the agent's stop-work age, DI-claim age, and PB-claim age. $\text{hist\_idx}$ is the last entry of \verb|n_a| and is treated as the experienceasset (\verb|vfoptions.experienceasset=1|). Section~\ref{sec:histidx} explains in detail.
\item \textbf{Exogenous state} \verb|n_z = 5|: health $h\in\{1,\ldots,5\}$, with age-dependent transition matrix \verb|pi_z_J| of shape $5\times 5\times N_j$.
\item $N_j = 56$, $\verb|agejshifter| = 24$.
\end{itemize}


\subsection{The \texttt{hist\_idx} experienceasset}\label{sec:histidx}

The single integer-valued state $\text{hist\_idx}\in\{1,\ldots,1237\}$ does the heavy lifting in this replication. It replaces what a more naive design would carry as four separate states: a DI-claim flag, a PB-claim flag, a stop-work-age tracker, and an Average Pension Points (AP) state. This section explains what $\text{hist\_idx}$ encodes, why each piece is needed to support the model's exercises, why we use a compact encoding rather than the unconstrained cartesian product, and how the integer is unpacked and updated.

\subsubsection{What it encodes}

$\text{hist\_idx}$ packs three ``employment-history coordinates'' into a single integer:
\begin{align*}
 \text{swa} &\;=\;\text{first stop-work age}\quad\text{(model period; 0 = never stopped),}\\
 \text{dia} &\;=\;\text{first DI-claim age}\quad\text{(model period; 0 = never claimed DI),}\\
 \text{pba} &\;=\;\text{first PB-claim age}\quad\text{(model period; 0 = never claimed PB).}
\end{align*}
All three are absorbing. Work is absorbing once stopped (consistent with LW2015's ``each possible employment history'' phrasing, p.129, and explicit in the LW2016 replication code's loop structure). DI is absorbing pre-65 and converts mechanically to PB at age 65 with the same benefit amount, so no separate ``DI-then-PB'' flag is needed. PB is absorbing once claimed.

\subsubsection{Why each coordinate is needed}

A reasonable question is whether all three coordinates are necessary. They are not all logically independent (the reachable region in Section~\ref{sec:reach} shows the dependencies), but conceptually each one is the minimal carrier of information for one of the model's benefit formulas:

\begin{description}
 \item[swa is needed for both OPB and the post-reform NDC pension.] The pre-reform blue-collar OPB depends on average pension points from real ages 55--59 (LW2015 p.130) --- the relevant slice of the wage history is pinned by swa. The pre-reform white-collar OPB depends on ``income at the time of retirement'' (p.130), i.e.\ wage at swa$-1$. Most importantly, the post-reform NDC pension recalculates the benefit when the agent stops working (LW2015 p.133: \emph{``The benefit is then recalculated when the individual stops working''}), so swa is the AP-finalization age in the post-reform system. Without swa as state, this rule cannot be implemented.
 \item[dia is needed for DI and the OPB DI top-up.] The DI benefit in both pre- and post-reform systems uses an AP projection that extends the ``contributing-years denominator'' to age 65 by filling in projected years at the rate of the last 3 pre-DI years (LW2015 p.131). That projection depends on \emph{when} DI was claimed, not just on the current age. The OPB DI top-up (LW2015 p.131; 15\% of pre-DI wage plus white-collar brackets) similarly depends on wage at dia$-1$, i.e.\ on dia.
 \item[pba is needed for the PB actuarial adjustment.] LW2015 p.130: $-0.5$\%-points per month early for the regular pension, $+0.7$\%/month delayed. That schedule is a function of the \emph{claim} age, which can differ from the stop-work age (the model allows claiming PB while still working).
\end{description}

These three are exactly the indices LW2016's hand-rolled code carries (\verb|i0, i1, i2| in lch1.m of the RED archive), and they are what the LW2015 paper's prose on p.129 (``\emph{for each possible employment history, consisting of disability, pension and retirement decisions}'') refers to.

\subsubsection{Reachable region}\label{sec:reach}

The unconstrained cartesian product of $(\text{swa}, \text{dia}, \text{pba})$ has $57 \times 41 \times 21 = 49{,}077$ entries:
\begin{itemize}
 \item swa ranges over $\{0,1,\ldots,56\}$ --- 56 model ages plus the sentinel 0 for ``never stopped''.
 \item dia ranges over $\{0, 1, \ldots, 40\}$ --- model ages 1--40 (real ages 25--64, the legal DI-claim window) plus 0 for ``never''.
 \item pba ranges over $\{0, 37, 38, \ldots, 56\}$ --- model ages 37--56 (real ages 61--80, the legal PB-claim window) plus 0 for ``never''.
\end{itemize}
But the model's structural constraints rule out the vast majority of this product. Two facts do the work:

\begin{enumerate}
 \item Work is absorbing once stopped, and DI implies stop-work (the agent cannot work on DI). So whenever dia $\ne 0$, we have $\text{swa} = \text{dia}$ --- one degree of freedom collapses.
 \item DI and PB are mutually exclusive pre-65, and the 65-conversion is mechanical (no separate PB-claim event). So on the DI branch, $\text{pba} = 0$ always.
\end{enumerate}

The reachable region therefore splits into two contiguous blocks:

\begin{itemize}
 \item \textbf{Block A (no DI ever):} $\text{dia} = 0$, with $(\text{swa}, \text{pba})$ varying freely over $\{0,\ldots,56\} \times \{0, 37, \ldots, 56\}$. That is $57 \times 21 = 1197$ tuples.
 \item \textbf{Block B (DI claimed):} $\text{dia} \in \{1,\ldots,40\}$, with $\text{swa} = \text{dia}$ (forced) and $\text{pba} = 0$. That is $40$ tuples.
\end{itemize}

Total reachable: $1197 + 40 = 1237$ tuples, about 2.5\% of the unconstrained product.

\subsubsection{The integer encoding}

We map the reachable region to integers $\{1,\ldots,1237\}$ by laying the two blocks out contiguously:

\paragraph{Block A (indices 1--1197).} Define $\text{pba\_idx} = 1$ if $\text{pba} = 0$, else $\text{pba\_idx} = \text{pba} - 35$ (giving $\text{pba\_idx}\in\{1,\ldots,21\}$). Then
\[
 \text{hist\_idx} \;=\; \text{swa} \cdot 21 + \text{pba\_idx}.
\]

\paragraph{Block B (indices 1198--1237).}
\[
 \text{hist\_idx} \;=\; 1197 + \text{dia}.
\]

\paragraph{Decoding.} Given $\text{hist\_idx}$:
\begin{itemize}
 \item If $\text{hist\_idx} \le 1197$:
 \begin{align*}
  \text{swa} &= \lfloor (\text{hist\_idx}-1)/21 \rfloor, \\
  \text{pba\_idx} &= ((\text{hist\_idx}-1) \bmod 21) + 1, \\
  \text{pba} &= \begin{cases} 0 & \text{if } \text{pba\_idx} = 1 \\ \text{pba\_idx} + 35 & \text{otherwise}\end{cases},\\
  \text{dia} &= 0.
 \end{align*}
 \item Otherwise (DI branch):
 \begin{align*}
  \text{dia} &= \text{hist\_idx} - 1197,\quad \text{swa} = \text{dia},\quad \text{pba} = 0.
 \end{align*}
\end{itemize}

The encode/decode arithmetic is duplicated inline in \verb|LuanWallenius2015_aprimeFn.m| and \verb|LuanWallenius2015_ReturnFn.m|. Factoring it into a helper file would be slightly cleaner but would prevent toolkit-side inlining --- both of those functions are on the inner loop of the value-function iteration and called millions of times.

\subsubsection{The decision variable d and the transition rule}

The decision $d \in \{1, 2, 3, 4, 5\}$ enumerates joint labor / claim actions:
\begin{itemize}
 \item $d=1$: keep working, no claim change. Feasible only if $\text{swa} = 0$ and the agent is not on DI.
 \item $d=2$: stop working (or stay stopped), no claim change. Always feasible; forced choice when on DI.
 \item $d=3$: apply for DI (forces stop work). Feasible only if $\text{dia} = 0$, $\text{agej} \le \text{agej\_di\_max}$, $\text{pba} = 0$, and $h \le 2$.
 \item $d=4$: claim PB and stop working. Feasible only if $\text{pba} = 0$, $\text{agej} \ge \text{agej\_pb\_first}$, and $\text{dia} = 0$.
 \item $d=5$: claim PB and keep working. Same as $d=4$ but also requires $\text{swa} = 0$.
\end{itemize}
Eligibility constraints are enforced inside the return function (\verb|LuanWallenius2015_ReturnFn.m|) by returning $-\infty$ for infeasible $(d, \text{state})$ pairs; the aprimeFn just executes the formal transition without checking eligibility, knowing that infeasible-action paths will be excluded by $-\infty$ utilities.

\paragraph{aprimeFn.} The history-state transition is a pure function of $(d, \text{hist\_idx}, \text{agej})$:
\begin{enumerate}
 \item Decode $\text{hist\_idx} \to (\text{swa}, \text{dia}, \text{pba})$.
 \item Apply the $d$-specific update rule:
\begin{align*}
 d=1: &\quad (\text{swa}', \text{dia}', \text{pba}') = (\text{swa}, \text{dia}, \text{pba}). \\
 d=2: &\quad \text{swa}' = \begin{cases}\text{agej} & \text{if swa}=0 \\ \text{swa} & \text{otherwise}\end{cases},\; \text{dia}'=\text{dia},\; \text{pba}'=\text{pba}.\\
 d=3: &\quad \text{if feasible: } (\text{swa}', \text{dia}', \text{pba}') = (\text{agej}, \text{agej}, \text{pba}). \\
 d=4: &\quad \text{if feasible: } \text{pba}' = \text{agej},\; \text{swa}' = \max(\text{swa}, \text{agej})|_{\text{first-stop}},\; \text{dia}'=\text{dia}. \\
 d=5: &\quad \text{if feasible: } \text{pba}' = \text{agej},\; \text{swa}' = \text{swa},\; \text{dia}'=\text{dia}.
\end{align*}
 (Where ``if feasible'' means the eligibility test in the action description above; otherwise the state is left unchanged --- the $-\infty$ utility in the return function makes this branch ignored by the argmax.)
 \item Encode $(\text{swa}', \text{dia}', \text{pba}') \to \text{hist\_idx}'$.
\end{enumerate}

\paragraph{Toolkit plumbing.} Because the transition depends only on the experienceasset, the decision, and the age --- not on $z$ or any other state --- this fits the plain experienceasset signature:
\begin{verbatim}
vfoptions.experienceasset = 1;
vfoptions.aprimeFn = @(d, hist_idx, agej, agej_pb_first, agej_di_max) ...
                     LuanWallenius2015_aprimeFn(d, hist_idx, agej, ...
                                                agej_pb_first, agej_di_max);
\end{verbatim}
Because the grid \verb|hist_grid = (1:1237)'| is integer-valued and the aprimeFn always returns an exact grid point, the toolkit's experienceasset lottery degenerates to a point mass at every transition --- no interpolation occurs, no precision is lost, and the policy function lives on the natural integer grid.

\subsubsection{Benefits as on-the-fly functions of \texttt{hist\_idx}}

All five benefit calculations (AP, PB, DI, OPB, OPB-DI-top-up) are pure functions of the decoded $(\text{swa}, \text{dia}, \text{pba})$ plus current age and PType-specific wage parameters. They are \emph{not} precomputed into matrices; instead, each lives in its own \verb|.m| file (\verb|LuanWallenius2015_AP.m|, \verb|_PB.m|, \verb|_DI.m|, \verb|_OPB.m|, \verb|_OPB_DI.m|) and is invoked by the return function. The AP function in particular implements the literal top-15-best rule from LW2015 p.130, with the 3-year-pre-DI projection that LW2015 p.131 describes for DI claimants (matching LW2016's \verb|penpoints.m|). The OPB and OPB-DI formulas implement the blue/white-collar piecewise rules from LW2015 p.130--131, including the 0.6\%/month early-claim reduction for white-collar and the 0.7\%/month late-claim bonus for blue-collar (LW2016 \verb|Toccl.m|, \verb|Tochs.m|, \verb|Tdioc.m|).

The runtime cost is an $O(N_j)$ pass through the wage profile inside each return-function call. A later optimization pass could lift each benefit into a precomputed lookup table indexed by $(\text{swa}, \text{dia}, \text{pba}, \text{PType}, \text{agej})$ without changing the public function signatures --- but for the current replication clarity wins over speed.

\subsubsection{Initial distribution}

All agents begin at $\text{hist\_idx} = 1$, which decodes to $\text{swa} = 0$, $\text{dia} = 0$, $\text{pba} = 0$ (``never claimed anything''). Combined with $k = 0$ and $h = 5$ (very good), this matches the paper's initial condition --- LW2015 p.129: \emph{``All agents start out with zero capital and in very good health.''}


\subsection{Health-transition Markov chain}\label{sec:health}

LW2015 Table~4 specifies three exogenous health shocks but does not aggregate them into an explicit one-step Markov transition. We work out the interpretation here.

\subsubsection{The three shocks}

Each period each agent draws:
\begin{itemize}
 \item a small negative shock, $\Delta h = -1$, with probability $p^l$;
 \item a large negative shock, $\Delta h = -3$, with probability $p^h$;
 \item a positive shock, $\Delta h = +1$, with probability $p^i$.
\end{itemize}
$p^l$ rises linearly in age from $0.1$ at age 25 to $0.6$ at age 80, with the LW2015 Table 4 override $p^l_1 = 0.6$ (constant) at $h=1$. $p^h$ is constant $0.01$, with override $p^h_1 = 0.1$ at $h=1$. $p^i$ varies with both age and health: it falls linearly from $\{0.6, 0.7, 0.8, 0.9, 1.0\}$ for $h \in \{1, 2, 3, 4, 5\}$ at age 25 to $0.5$ at age 80.

\subsubsection{The aggregation ambiguity}

The paper does not say how the three shock probabilities combine into a Markov transition $\pi_z(h \to h')$. The natural reading is ``one of the three events fires, residual $1 - p^l - p^h - p^i$ is no change''. But at several $(h, \text{age})$ cells the three probabilities sum to more than 1:
\begin{center}
\begin{tabular}{lccc}
 & $p^l$ & $p^h$ & $p^i$ \\
\hline
young $h=4$ & $0.10$ & $0.01$ & $0.90$ \\
old $h=4$   & $0.60$ & $0.01$ & $0.50$ \\
old $h=3$   & $0.60$ & $0.01$ & $0.50$
\end{tabular}
\end{center}
At each of these the three numbers cannot simultaneously be transition probabilities of mutually exclusive events. The mutually-exclusive reading therefore requires ad-hoc clipping or renormalization, both of which throw away ``stay'' mass at exactly the ages and health states most relevant for the paper's results.

\subsubsection{Reading A --- independent Bernoulli draws}

We adopt instead the reading that the three shocks fire \emph{independently}. Each period each shock fires (Bernoulli) with its Table-4 probability; the realized $\Delta h$ is the algebraic sum of the firing shock values; the new health state is the cap $h' = \max(1, \min(5, h + \Delta h))$.

The joint distribution has $2^3 = 8$ outcomes. For a given $(h, p^l, p^h, p^i)$:
\begin{center}
\begin{tabular}{lccc}
event combination & probability & $\Delta h$ & $h'$ \\
\hline
none                 & $(1-p^l)(1-p^h)(1-p^i)$           & $0$  & $h$ \\
only $p^l$           & $p^l(1-p^h)(1-p^i)$               & $-1$ & $\max(h-1, 1)$ \\
only $p^h$           & $(1-p^l)p^h(1-p^i)$               & $-3$ & $\max(h-3, 1)$ \\
only $p^i$           & $(1-p^l)(1-p^h)p^i$               & $+1$ & $\min(h+1, 5)$ \\
$p^l$ and $p^h$      & $p^l p^h (1-p^i)$                 & $-4$ & $1$ \\
$p^l$ and $p^i$      & $p^l (1-p^h) p^i$                 & $0$  & $h$ \\
$p^h$ and $p^i$      & $(1-p^l) p^h p^i$                 & $-2$ & $\max(h-2, 1)$ \\
all three            & $p^l p^h p^i$                     & $-3$ & $\max(h-3, 1)$
\end{tabular}
\end{center}
The transition $\pi_z(h \to h')$ is the sum of all event probabilities whose net change lands at $h'$. The resulting matrix is denser than under the mutually-exclusive reading --- every $h \to h'$ within four-step reach has positive mass, not just the three the table's headline numbers suggest --- but it requires no clipping and uses every published probability at face value.

\subsubsection{Qualitative consequences}

The independent-draw structure gives the positive shock a partial ``shield'' role. At $h=5$ young, $p^i = 1.0$ means the positive shock fires every period with certainty. A coincident small negative shock then nets to $0$ (stay at $h=5$); a coincident large negative shock nets to $-2$ (lands at $h=3$ instead of $h=2$). At unhealthy states the same mechanism in reverse: at $h=1$ the elevated $p^l = 0.6$ doesn't move the agent through the cap, but it does eat into the positive-shock recovery rate because the joint event ``positive and small negative'' nets to $0$ (stay at $h=1$). Both effects produce health persistence at the boundary states, which the paper notes as a target on p.128.

\subsubsection{Why Reading A}

Reading A uses every Table-4 parameter at face value, requires no renormalization, has a clean derivation, and is the closest match to LW2016's stylistic preference for independent-event products: in [Compusswe/Version 1/lch1.m] the negative shocks $(p^l, p^h)$ and the (endogenous) investment-success probability $p_e$ are combined as $(1-p^l-p^h) \cdot p_e + \ldots$ --- a literal joint-probability expansion of independent draws.

\subsubsection{Empirical validation}

The reading-choice question is settled directly against the paper's own reported model output: since health is purely exogenous in LW2015, the unconditional health distribution at any age is just the initial $P(h=5)=1$ propagated through $\pi_z$, independent of all other model parameters. LW2015 p.132 reports their model's distribution at real age 70 explicitly, and propagating under each reading gives an L1 distance to those numbers of $\mathbf{5.6}$ percentage points under Reading A vs $\mathbf{34.4}$ under Reading B. Reading A is decisively confirmed and Reading B ruled out. Appendix~\ref{app:health-test} reproduces the full propagation calculation, the test code, and the result table.


\subsection{Calibration shortcuts}

What we still depart from in the paper:

\begin{itemize}
\item \textbf{Wage profiles.} The four life-cycle earnings profiles in LW2015 Fig.~1 (LINDA 1968--1997) are estimated by regression on age, age$^2$, and individual fixed effects, but the regression coefficients are not printed in the paper, the SSE working-paper version (WP 741), or any online appendix or replication package we can locate. We \emph{cannot} borrow from the LW2016 sequel replication code (RED 14-161) either, because LW2016 publishes only \emph{two} regressed wage profiles in \verb|Compusswe/Version 1/w.m| (high school and college) --- the within-education fixed-effect dimension that LW2015 uses for the NCL/NCH and CL/CH split was dropped in LW2016. Pending either the original LW2015 LINDA regression output or LINDA access to re-estimate, we eyeball four hump-shaped quadratics from Fig.~1 with peak at real age 50; see the discussion in Section~\ref{sec:wage-profiles} for details.
\item \textbf{Lump-sum transfer in the pre-reform baseline.} Held at $T = 0.4\,\text{BA}$ (the paper's reported equilibrium value; $\approx 15{,}000$ SEK / 36{,}300 SEK $\approx 0.4\,\text{BA}$) rather than re-derived via a budget-balance loop. For the post-reform regimes $T$ is solved as a GE price via the toolkit's \texttt{HeteroAgentStationaryEqm\_Case1\_FHorz\_PType} with an inline \texttt{GeneralEqmEqns.budget = @(surplus, T) surplus - T}; see \texttt{doPart(3)} (regular-only GE) and \texttt{doPart(5)} (full reform GE) in the driver.
\item \textbf{Disutility-from-work and health-shock probabilities.} Set to LW2015 Table 4 values; not yet jointly recalibrated to match the employment and DI-incidence age profiles (LW2015 Tables 1--2).
\item \textbf{PType mass.} Set to LW2015's 17\% college / 83\% non-college sample share, split 50/50 low/high within each education group: $(0.415, 0.415, 0.085, 0.085)$.
\end{itemize}

What is now done at full paper fidelity (where the earlier draft of this replication had simplified):
\begin{itemize}
\item AP is the literal top-15-best rule with 3-year-pre-DI projection.
\item PB benefit includes the $\min(\text{yrs}/30,1)$ proration and the $\pm 0.5/0.7\%$/month actuarial adjustment.
\item DI benefit includes the projection-to-65 ``assumed pension points'' rule.
\item OPB supports early claim from age 55 for white-collar with the 0.6\%/month reduction, and late-claim bonus for blue-collar to age 70.
\item OPB DI top-up paid alongside DI (15\% of pre-DI wage for blue-collar; white-collar piecewise).
\end{itemize}

The next step in this replication would be (i) regressing the wage profiles from a comparable Swedish micro-dataset, and (ii) jointly calibrating the disutility levels $b(h)$ and the health-shock probabilities to match the employment and DI-incidence age profiles in LW2015 Tables 1--2 (the targets the paper itself calibrates to).
